%!TeX program = lualatex
\renewcommand{\thelektion}{L05 \textperiodcentered{} ASCII \& UTF-8}

\begin{frame}\lessontitle{Lektion 05}{Kodierung von Buchstaben}{ASCII und UTF-8 \textperiodcentered{} Skript, Kapitel 3}\end{frame}

\begin{frame}\FDUtitle{Wie kodiert man Buchstaben?}
\begin{center}
Computer speichern nur Nullen und Einsen. Wie könnte man damit \textbf{Buchstaben} repräsentieren?
\end{center}
\end{frame}

\begin{frame}\FDUtitle{ASCII}
\gls{ASCII} -- eine der ersten Text-Kodierungen
\begin{itemize}\setlength\itemsep{6pt}
	\item 1 Byte pro Zeichen -- z.\,B.\ \texttt{0100 0001} $\rightarrow$ A
	\item 1 Bit war für Fehlerkontrolle reserviert (Paritätsbit)
\end{itemize}
\only<2->{
\begin{center}
	\textbf{Frage}: Wie viele Zeichen kann man mit einem Byte abbilden? Und mit nur 7 nutzbaren Bits (ASCII)?
\end{center}
}
\only<3->{\begin{center}\Large $2^8 = 256,\quad 2^7 = 128$\end{center}}
\end{frame}

\begin{frame}[fragile]\FDUtitle{ASCII-Tabelle (Auszug)}
\begin{table}[H]
	\centering
	\tiny
	\begin{tabular}{cccc}
		\toprule
		Dezimal & Hexadezimal & Binär & Zeichen \\
		\midrule
		32 & 20 & 00100000 & Space \\
		48 & 30 & 00110000 & 0 \\
		49 & 31 & 00110001 & 1 \\\midrule
		... & ... & ... & ... \\\midrule
		65 & 41 & 01000001 & A \\
		66 & 42 & 01000010 & B \\
		67 & 43 & 01000011 & C \\\midrule
		... & ... & ... & ... \\\midrule
		97 & 61 & 01100001 & a \\
		98 & 62 & 01100010 & b \\
		\bottomrule
	\end{tabular}
\end{table}
\begin{center}Vollständige Tabelle: Skript, Kapitel 3.\end{center}
\end{frame}

\begin{frame}[fragile]\FDUtitle{Entschlüsseln Sie!}
\begin{center}
\texttt{01101000 01100101 01101100 01101100 01101111}
\end{center}
\begin{center}
Nutzen Sie die ASCII-Tabelle im Skript: welches Wort ist das?
\end{center}
\only<2->{\begin{center}\Large \texttt{hello}\end{center}}
\end{frame}

\begin{frame}\FDUtitle{Die Grenzen von ASCII}
\begin{center}
128 Zeichen reichen bei weitem nicht für Umlaute, Akzente, Emojis, Sprachen wie Chinesisch, Arabisch, \ldots
\end{center}
\vspace{4mm}
\begin{center}
$\rightarrow$ \gls{UTF-8}: bis zu \textbf{4 Bytes} pro Zeichen
\end{center}
\only<2->{
\begin{center}\textbf{Frage}: Wie viele Zeichen kann man mit 4 Bytes maximal abbilden?\end{center}
}
\only<3->{\begin{center}\Large $2^{32} = 4\,294\,967\,296$ (etwas weniger nutzbar, da Bits zur Längenangabe gebraucht werden)\end{center}}
\end{frame}

\begin{frame}{Auftrag}{Kapitel 3: ASCII \& UTF-8}
\aufgabebox{\textbf{Aufgabe 3.3}: Wie viele Zeichen lassen sich mit 7-Bit ASCII abbilden?}\\[3mm]
\aufgabebox{\textbf{Aufgabe 3.4}: Grösste binäre Zahl mit 8 Bits und in Hexadezimal}\\[3mm]
\aufgabebox{\textbf{Aufgabe 3.5}: Maximale Anzahl Zeichen bei 4-Byte UTF-8}
\end{frame}
